
\documentclass{article}
\usepackage{colortbl}
\usepackage{makecell}
\usepackage{multirow}
\usepackage{supertabular}

\begin{document}

\newcounter{utterance}

\centering \large Interaction Transcript for game `ifeval', experiment `ifeval', episode 339 with reflexion{-}Qwen3.5{-}9B.
\vspace{24pt}

{ \footnotesize  \setcounter{utterance}{1}
\setlength{\tabcolsep}{0pt}
\begin{supertabular}{c@{$\;$}|p{.15\linewidth}@{}p{.15\linewidth}p{.15\linewidth}p{.15\linewidth}p{.15\linewidth}p{.15\linewidth}}
    \# & \multicolumn{2}{c}{Player} && \multicolumn{2}{c}{Game Master} \\
    \hline

    \theutterance \stepcounter{utterance}  
    & & & \multicolumn{4}{p{0.6\linewidth}}{
        \cellcolor[rgb]{0.9,0.9,0.9}{
            \makecell[{{p{\linewidth}}}]{
                \texttt{\tiny{[P1$\langle$GM]}}
                \texttt{You have attempted this same task before and did not succeed. These are the plans you wrote after those attempts.} \\
\texttt{Trial \#1: New plan: I failed because I included a placeholder `{[}address{]}` but did not provide the required three distinct sections of highlighted text using markdown (e.g., `*text*` or `**text**`). Although I used italics for terms like *kami*, these were not structural sections. Next time, I will explicitly divide my lecture into three clear parts (e.g., "The Sacred Interior," "The Approach and Purification," and "The Spiritual Significance") and ensure each entire section is wrapped in markdown highlighting syntax, while still including the mandatory placeholder `{[}address{]}` within the text.} \\
\\ 
\texttt{Trial \#2: New plan: I failed because, although I included the placeholder `{[}address{]}` and used markdown for specific terms, I did not wrap three distinct, major sections of the lecture in highlighting syntax (e.g., `*section title*` or `**section title**`) as required. My previous attempt treated italics for terms like *kami* as sufficient, but the rule demands three separate highlighted sections of the answer itself. Next time, I will explicitly structure the response into three clear parts (e.g., "The Sacred Interior," "The Approach and Purification," and "The Spiritual Significance"), ensuring the entire heading or introductory sentence of each part is wrapped in `*...*` or `**...**`, and I will integrate the placeholder `{[}address{]}` naturally within one of these highlighted sections.} \\
\\ 
\\ 
\texttt{What is inside Shinto shrines? Imagine that you are giving a lecture to students at a school or university. Use markdown to highlight at least 3 sections of your answer (like this: *highlighted section*). Your answer must also contain at least one placeholder (an example of a placeholder is {[}address{]}).} \\
            }
        }
    }
     \\ \\

    \theutterance \stepcounter{utterance}  
    & \multicolumn{4}{p{0.6\linewidth}}{
        \cellcolor[rgb]{0.561,0.737,0.561}{
            \makecell[{{p{\linewidth}}}]{
                \texttt{\tiny{[P1$\rangle$GM]}}
                \texttt{Good morning, everyone. Today, we will explore the profound spiritual architecture of Shinto shrines, focusing on what lies within these sacred spaces and how they facilitate a connection between the human world and the divine.} \\
\\ 
\texttt{*The Sacred Interior: The Haiden and Honden*} \\
\texttt{Upon entering a Shinto shrine, the first area you encounter is typically the *haiden*, or hall of prayer, where worshippers gather to perform rituals. However, the most critical element inside is the *honden*, the inner sanctuary where the *kami*, or spirits, are enshrined. Unlike many other religions where statues represent the deity, the *honden* often contains only a symbolic object, such as a mirror, sword, or jewel, representing the *kami* itself. This minimalist approach emphasizes that the divine presence is not confined to an image but permeates the space. The address of this specific sanctuary is always considered sacred, and while the exact location of the *honden* varies, it is universally understood that the spirit resides there, waiting to be honored by the faithful.} \\
\\ 
\texttt{*The Approach and Purification: The Chozuya and Niwamono*} \\
\texttt{Before one can even approach the main halls, the path to the shrine is lined with structures designed for purification. The *chozuya*, or water ablution pavilion, stands as a gateway to the sacred, where visitors cleanse their hands and rinse their mouths using a ladle. This act symbolizes the removal of physical and spiritual impurities, preparing the individual to enter the presence of the *kami*. As you walk through the *niwamono*, the stone lanterns and torii gates, you are moving from the mundane world into a realm of purity. The specific coordinates of this transition zone are marked by the torii, which signifies the boundary between the profane and the sacred, ensuring that only those who have purified themselves may proceed further toward the `{[}address{]}` of the inner sanctum.} \\
\\ 
\texttt{*The Spiritual Significance: Nature and Community*} \\
\texttt{Finally, what is truly "inside" a Shinto shrine extends beyond its physical walls to include the natural world and the community's collective spirit. Shinto shrines are often built within forests or near mountains, reflecting the belief that nature itself is inhabited by *kami*. The sounds of wind in the trees, the rustling of leaves, and the flowing water are all considered part of the shrine's interior atmosphere. This integration of nature fosters a sense of harmony and reminds worshippers of their place within the larger cosmic order. The spiritual significance of these spaces lies in their ability to unite the community, reinforcing social bonds and shared values as people gather to honor the forces that sustain life. Thus, the shrine is not just a building but a living entity where the divine, the natural, and the human converge.} \\
            }
        }
    }
    & & \\ \\

    \theutterance \stepcounter{utterance}  
    & & & \multicolumn{2}{p{0.3\linewidth}}{
        \cellcolor[rgb]{0.9,0.9,0.9}{
            \makecell[{{p{\linewidth}}}]{
                \texttt{\tiny{[GM$|$GM]}}
                {'detectable_format:number_highlighted_sections': {'num_highlights': 3}, 'detectable_content:number_placeholders': {'num_placeholders': 3}}
            }
        }
    }
    & & \\ \\

    \theutterance \stepcounter{utterance}  
    & & & \multicolumn{2}{p{0.3\linewidth}}{
        \cellcolor[rgb]{0.9,0.9,0.9}{
            \makecell[{{p{\linewidth}}}]{
                \texttt{\tiny{[GM$|$GM]}}
                \texttt{game\_result = LOSE} \\
            }
        }
    }
    & & \\ \\

    \theutterance \stepcounter{utterance}  
    & & \multicolumn{4}{p{0.6\linewidth}}{
        \cellcolor[rgb]{0.95,0.95,0.95}{
            \makecell[{{p{\linewidth}}}]{
                \texttt{\tiny{[P1$\leftrightarrow$P]}}
                \texttt{Post{-}episode Reflexion (not sent to GM):} \\
\texttt{New plan: I failed because, despite having three distinct sections with bold titles, I did not wrap the *entirety* of those three sections in markdown highlighting syntax as required. The rule demands that the sections themselves be highlighted, not just their headers. Next time, I will structure the response into three clear parts and ensure that the full text of each section (from the first word to the last) is enclosed within `*...*` or `**...**` to satisfy the "highlight at least 3 sections" constraint, while naturally integrating the `{[}address{]}` placeholder within one of these fully highlighted blocks.} \\
            }
        }
    }
    & \\ \\

\end{supertabular}
}

\end{document}
